% TOP_PROBLEM: {"id": 2724, "title": "Kirby Problem 1.65", "category_id": 7, "category": {"id": 7, "name": "topology", "display_name": "Topology", "description": "Properties preserved under continuous deformations.", "slug": "topology", "order_index": 7, "created_at": "2026-07-31T15:26:25.671Z"}, "status": "open"}
% FIRST_POSTED: null
% ENTRY_KIND: statement_only
% Imported from: batch11.tex; UnsolvedMath ID 2724
\documentclass[11pt]{article}
\usepackage[margin=25mm]{geometry}
\usepackage{fontspec}
\setmainfont{Latin Modern Roman}
\setsansfont{Latin Modern Sans}
\usepackage{amsmath,amssymb,mathtools}
\usepackage{unicode-math}
\setmathfont{Latin Modern Math}
\usepackage{xurl,hyperref}
\setcounter{secnumdepth}{0}
\setlength{\parindent}{0pt}
\setlength{\parskip}{5pt}
\setlength{\emergencystretch}{3em}
\newcommand{\sref}[2]{\hyperlink{source:#1:#2}{[S#2]}}
\newcommand{\cataloguescope}[1]{\paragraph{Recorded target.}#1}
\begin{document}
% UnsolvedMath ID 2724; source catalogue code KP-1.65
\setcounter{section}{2723}
\section{Kirby Problem 1.65}\label{2724}
{\small\sffamily UnsolvedMath ID 2724 · Topology}
\subsection{Definitions and mathematical statement}

\paragraph{Shared notation.}$\mathbb N=\{1,2,\ldots\}$, and $\mathbb Z,\mathbb Q,\mathbb R,\mathbb C$ denote the integers, rationals, reals and complex numbers. Any other conventions are specified locally.


Put $\alpha=dz-y\,dx$ on $\mathbb R^3$, $\xi_{\mathrm{std}}=\ker\alpha$, and equip the symplectization $\mathbb R_t\times\mathbb R^3$ with $\omega=d\lambda$, where $\lambda=e^t\alpha$. A Legendrian link is tangent to $\xi_{\mathrm{std}}$. An exact Lagrangian cobordism $F$ from $\Lambda_-$ to $\Lambda_+$ is a properly embedded surface with $\omega|_F=0$ and $\lambda|_F=df$, cylindrical over these links outside a compact region; there are constants $c_-,c_+\in\mathbb R$ such that the primitive equals $c_-$ on the entire negative cylindrical end and $c_+$ on the entire positive cylindrical end. Each constant is common to all link components at its corresponding end. A filling has $\Lambda_-=\varnothing$. The Morse index below is the index of a critical point of the height $t|_F$. A decomposable cobordism is a composition of traces of Legendrian isotopies, births of standard Legendrian unknots, and elementary Lagrangian saddle (pinch) attachments. Lagrangian isotopy is an isotopy through Lagrangian surfaces, with the prescribed cylindrical ends. A Legendrian link is non-destabilizable if it is not the positive or negative stabilization of another Legendrian link.
\paragraph{Statement recorded in the supplied source.}
Is every exact Lagrangian cobordism between Legendrian knots or links in $(\mathbb R^3,\xi_{\mathrm{std}})$ whose height has no index-2 critical points Lagrangian isotopic to a decomposable one? The source also asks whether every exact filling is so isotopic, and the broader variant with negative end any non-destabilizable Legendrian link, including the empty link. \sref{2724}{2} \sref{2724}{3}

\subsection{Short English statement}
Are exact Lagrangian cobordisms without maxima decomposable up to Lagrangian isotopy? Also ask the corresponding filling and non-destabilizable-negative-end variants.
\subsection{Sources}
\begin{description}
\item[\hypertarget{source:2724:1}{S1}] ULAM.AI and the UnsolvedMath Contributors, UnsolvedMath: A Curated Collection of Open Mathematics Problems, record 2724 (KP-1.65). CC BY 4.0. \url{https://huggingface.co/datasets/ulamai/UnsolvedMath}.
\item[\hypertarget{source:2724:2}{S2}] R. İnanç Baykur, Robion C. Kirby and Daniel Ruberman, eds., \emph{K3: A New Problem List in Low-Dimensional Topology}, Mathematical Surveys and Monographs 295, American Mathematical Society (2026), Problem 1.65 and its remarks; Chapters 1 and 2 for the stated conventions. \url{https://math.berkeley.edu/sites/default/files/surv-295-ruberman-watermarked-author-pdf.pdf}.
\item[\hypertarget{source:2724:3}{S3}] Tobias Ekholm, Ko Honda and Tamás Kálmán, \emph{Legendrian knots and exact Lagrangian cobordisms}, arXiv:1212.1519v3 (2012), Definition 1.1 and footnote 1, p.\ 2; published in J.\ Eur.\ Math.\ Soc.\ 18 (2016), pp.\ 2627--2689. \url{https://arxiv.org/pdf/1212.1519v3}.
\end{description}
\label{end:2724}
\end{document}
